\section{A fixed-multitape comparison}
\label{sec:multitape}

The single selected block in Lemma~\ref{lem:visits} is essential to
the factored round rule. For example, a joint transition inspecting
two independently supplied bits may produce a specified outgoing
controller precisely when the bits are equal. The successful pairs
\(\{00,11\}\) cannot be a product of predicates on the individual
bits. An initial visit supplies this obstruction even with the full
earlier short history fixed.
This observation concerns this factorization, not the possibility
of other multitape simulation methods.

For completeness, the ordinary-rule part of
Lemma~\ref{lem:weighted-evaluation} gives the following comparison.
Here the access condition of Section~\ref{sec:introduction} has a tag
for each initial writable tape and each read-only head, with coordinates
measured from their fixed origins.

\begin{proposition}
\label{prop:multitape}
Fix a deterministic machine with a fixed number of ordinary writable
tapes and read-only input heads, unit movement and fixed origins for
all heads, cap-independent symbol functions, and a fixed accessor.
For every public input \(x\) and supplied binary cap \(T\ge2\) satisfying
the access condition, its finite-control and halting outcome admits a
deterministic simulation
in \(\softO(\sqrt T)\) work-space bits, excluding read-only input storage.
Simulation time is unrestricted.
\end{proposition}

\begin{proof}
If there are no writable tapes, a direct capped simulation stores the
finite-control state, elapsed time, terminal flags, and the fixed number
of read-only head displacements. Unit movement keeps every query within
the cap radius, so the controller and accessor use \(\polylog T\) work bits.
Hence assume the writable-tape count \(m\) is positive, and let
\(b=\lceil\sqrt T\rceil\). On each writable tape, use the represented
blocks of Lemma~\ref{lem:visits}: they contain the real symbols within
distance \(T\) of that tape's origin and a fixed filler outside.
All block words below refer to this represented computation.
The confinement argument applies to every head of every accepted
candidate controller. From elapsed time \(\tau\) and displacement at
most \(\tau\), after \(s\le T-\tau\) further unit steps its displacement
is at most \(\tau+s\le T\). Thus a joint visit never reads or writes an
exterior writable cell, and all read-only queries satisfy the access
condition. Invalid or terminal controllers make no underlying tape query.
The represented and real capped controller trajectories therefore agree.

Apply a common trial offset to all tapes.
The sum of crossing counts, averaged over offsets, is at most
\(mT/b\). End a joint visit when any writable head crosses, or on
halt/timeout. Some offset therefore needs at most
\(R=\lceil1+mT/b\rceil\) visits. A joint visit reads and updates
only its \(m\) selected blocks and is computable in
\(O((mb+1)\polylog T)\) space on those explicit blocks, including
public-symbol access.

For each offset, enumerate a complete list of \(R+1\) raw fixed-width
controller strings, including the prescribed initial one. Hold this list
in one global \(O(R\log(T+2))\)-bit buffer.
For a candidate list \(z_0,\ldots,z_R\), let \(Q_{t,h}\) denote
the block word produced after visit \(t\) for the block on tape \(h\) selected
by the incoming controller \(z_{t-1}\).
Its incoming block is the output \(Q_{s,h}\) from the last earlier
visit \(s<t\) selecting that same tape/block, or the initial block
if none exists. Each \(Q_{t,h}\), and each gate computing or checking
the outgoing controller, depends on at most \(m\) incoming block
words and the prescribed controller \(z_{t-1}\).
The last-update index is found by scanning the retained list.
All dependencies decrease visit number, so their depth is \(O(R)\),
even for malformed lists.

Use the same total parsing convention as in Lemma~\ref{lem:visits},
with a designated valid block on each tape for an invalid controller.
Each transition check compares its deterministic raw outgoing controller
bit for bit with the proposed string \(z_t\).
Apply the ordinary-rule evaluator with word width \(O(b)\), fixed
fan-in and path weight \(O(R)\).
Evaluate the transition checks one at a time. Induction from the
initial controller proves that a passing list is exactly the actual
capped computation: the last-update rule supplies each true
incoming block, and the next check forces the true outgoing
controller. Require a terminal last controller; disregard
insufficient offsets. A completed actual list exists for a good
offset, with stationary padding after termination.

The controller list is stored once, not copied into recursive
frames. Evaluator storage and that buffer together occupy
\(\softO(b+R)=\softO(\sqrt T)\).
Only the maximum storage over offset/list trials is needed.
\end{proof}

Proposition~\ref{prop:multitape} matches the known square-root
exponent but does not assert the sharper logarithmic factor of
Williams's Theorem~1.1~\cite{Williams2025}.
For a fixed ordinary multitape decider halting within
\(T\ge\max\{n,2\}\) steps on an input of length \(n\), the standard
one-writable-tape conversion halts with the same decision within
\(O(T^2)\) steps. Applying the form of Theorem~\ref{thm:main} without
a supplied running time to that converted decider would give
\(\softO(T^{4/5})\) space in the original time parameter.
This is a running-time comparison for halting deciders.
Thus neither the conversion nor the factorization argument
establishes a new general-multitape exponent.
